\documentclass[12pt,letterpaper]{report}
\usepackage{graphicx}
\usepackage{scrextend}
\usepackage{vmargin}
\usepackage{graphicx}
\usepackage{multirow}
\usepackage[utf8]{inputenc}
\usepackage[spanish]{babel}
\usepackage{multicol}
\usepackage{enumerate}
\usepackage{float}
\usepackage{amsmath, amsthm, amssymb, amsfonts}
\usepackage[usenames]{color}
\parindent=0mm
\pagestyle{empty}
\definecolor{miorange}{rgb}{0.91, 0.43, 0.0}
\begin{document}
\setmargins{2.5cm}      
{1.5cm}                     
{2cm}  
{24cm}                    
{10pt}                          
{1cm}                          
{0pt}                             
{2cm}
\begin{flushright}
\today
\end{flushright}
\vspace{-1.75cm}
\section*{Tarea Clase 11}
\subsection*{Obtenga las soluciones de las siguientes ecuaciones:}
\begin{enumerate}
\item Se quiere colocar un cable desde la cima de una torre de 25 metros altura hasta un punto situado a 50 metros de la base la torre. ¿Cuánto debe medir el cable?\\
La altura del triangulo es de 25, la base es de 50, usando el teorema de pitagoras encontraremos el otro lado, por lo que:
\begin{align*}
    z^2&=50^2+25^2\\
    z^2&=5(25^2)\\
    z&=25\sqrt{5}
\end{align*}
\item Una parcela de terreno cuadrado dispone de un camino de longitud $2\sqrt{2}$ kilómetros (segmento discontinuo) que la atraviesa según se muestra en la siguiente imagen. Calcular el área total de la parcela.\\
No viene la imagen, pero es la diagonal de un cuadrado, por lo que se tendria un triangulo con base y altura iguales, por lo que:
\begin{align*}
    (2\sqrt{2})^2 &=x^2 + x^2\\
    8&=2x^2\\
    x^2=4\\
    x=2
\end{align*}
por lo tanto el área es de 4
\item La hipotenusa de un triángulo rectángulo mide 10 metros y sus catetos miden x y x+2. ¿Cuánto miden los catetos?.\\
Desarrollando el teotema de pitagoras:
\begin{align*}
    x^2+(x+2)^2&=100\\
    x^2+x^2+4x+4&=100\\
    2x^2+4x-96&=0\\
    x^2+2x-48&=0\\
    (x-6)(x+8)&=0\\
    x=6,-8
\end{align*}
usando x=6, encontramos que el otro lado es 8
\item Se desea pintar una cuadrado inscrito en una circunferencia de radio 
$R=3$ como se muestra en la figura. Calcular el área del cuadrado.\\
Tampoco viene la figura pero este cuadrado es dentro del circulo, por lo que notaremos que el radio del circulo es la mitad de un lado del cuadrado, por ende un lado del cuadrado es igual a 6, por lo que su área es 36
\item Calcular cuánto mide el cateto b de un triángulo rectángulo si su otro cateto, a, y su hipotenusa, h, miden 
\begin{eqnarray*}
a=\frac{3\sqrt{2}}{2} \\
h=\frac{\sqrt{194}}{6}
\end{eqnarray*}
\item Hallar las medidas de los lados de una vela con forma de triángulo rectángulo si se quiere que tenga un área de 30 metros al cuadrado y que uno de sus catetos mida 5 metros para que se pueda colocar en el mástil.
Usando el teorema de pitagoras encontramos que:
\begin{align*}
    \left(\frac{\sqrt{194}}{6} \right)^2& =\left(\frac{3\sqrt{2}}{2} \right)^2 +x^2\\
    \frac{194}{36} &= \frac{18}{4} + x^2  \\
    x^2&= \frac{8}{9}\\
    x&= \frac{2\sqrt{2}}{3}
\end{align*}
\item Si el cateto de un triángulo rectángulo mide x y el otro mide el doble, obtener una fórmula para calcular la longitud de la hipotenusa en función del cateto menor, x.Utilizar la fórmula obtenida para calcular la hipotenusa cuando $x=\sqrt{5}$ y $x=2\sqrt{5}.$\\
Se tiene que a=x y b=2x, por lo que:
\begin{align*}
    c^2&=a^2+b^2\\
    c^2&=x^2+4x^2\\
    c^2&=5x^2\\
    c&=x\sqrt{5}
\end{align*}
sustituyendo los datos que nos dan:
\begin{align*}
    c&=(\sqrt{5})\sqrt{5}\\
    c&=5
\end{align*}
\end{enumerate}
\end{document}